\section{Server API}
\label{sec:ServerApi}

This chapter documents the WebSocket and REST API endpoints exposed by the
Archimesh server.

\subsection{WebSocket endpoint}

\begin{description}
\item[Path] \texttt{/models/\{modelId\}/stream}
\item[Protocol] WebSocket (RFC 6455)
\item[Sub-protocol] None (plain JSON text frames)
\end{description}

The WebSocket endpoint handles real-time collaboration. Each connection is
scoped to one model.

\subsubsection{WebSocket lifecycle}

\begin{enumerate}
\item Client opens a WebSocket connection to the endpoint.
\item Server performs authorization checks (see
  Section~\ref{sec:Authorization}) and either accepts the connection or
  closes it with policy violation status.
\item Client sends a \texttt{Join} message to request model state.
\item Server responds with \texttt{CheckoutSnapshot} or
  \texttt{CheckoutDelta}.
\item Client sends \texttt{SubmitOps}, \texttt{AcquireLock},
  \texttt{ReleaseLock} and \texttt{Presence} messages during the session.
\item Server broadcasts \texttt{OpsAccepted}, \texttt{OpsBroadcast},
  \texttt{LockEvent}, \texttt{PresenceBroadcast} and \texttt{Error}
  messages to connected sessions.
\item Client or server closes the WebSocket. On close the server unregisters
  the session and releases held locks.
\end{enumerate}

\subsubsection{Inbound messages (client to server)}

All inbound messages are wrapped in a \texttt{ClientEnvelope}:

\begin{lstlisting}[caption={ClientEnvelope shape}]
{
  "type": "<MessageType>",
  "payload": { ... }
}
\end{lstlisting}

\begin{longtable}{p{32mm}p{52mm}p{70mm}}
\toprule
Type & Purpose & Key payload fields \\
\midrule
\texttt{Join} & Request model checkout &
  \texttt{lastSeenRevision} (optional),
  \texttt{ref} (optional; \texttt{HEAD} or tag name),
  \texttt{actor} (\texttt{userId}, \texttt{sessionId}) \\
\texttt{SubmitOps} & Submit one batch of operations &
  \texttt{baseRevision},
  \texttt{opBatchId},
  \texttt{actor},
  \texttt{ops} (array) \\
\texttt{AcquireLock} & Acquire a notation lease &
  \texttt{actor},
  \texttt{targets} (array of entity ids),
  \texttt{ttlMs} \\
\texttt{ReleaseLock} & Release a notation lease &
  \texttt{actor},
  \texttt{targets} (array of entity ids) \\
\texttt{Presence} & Broadcast cursor/selection &
  \texttt{actor},
  \texttt{viewId},
  \texttt{selection} (array),
  \texttt{cursor} (object) \\
\bottomrule
\end{longtable}

\subsubsection{Join and checkout}

When the server receives a \texttt{Join}:

\begin{enumerate}
\item The model must exist (registered via admin API).
\item Authorization is checked (join action).
\item The requested ref is resolved (\texttt{HEAD} or tag name).
\item If \texttt{lastSeenRevision} is present and a delta is available
  (the revision gap is within the op-log window), the server responds
  with \texttt{CheckoutDelta} containing the missing operations.
\item Otherwise, the server responds with \texttt{CheckoutSnapshot}
  containing the full materialized model state at the requested ref.
\end{enumerate}

\subsubsection{SubmitOps processing}

The \texttt{SubmitOps} pipeline (implemented in \texttt{ArchimeshService}):

\begin{enumerate}
\item Model registration and writable ref checks.
\item Validation of the op batch structure.
\item Normalization and expansion (bootstrap view ref normalization,
  cascade delete expansion).
\item Revision-ahead check: if \texttt{baseRevision > headRevision},
  broadcast \texttt{Error(REVISION\_AHEAD)} to all sessions.
\item Precondition validation: all entity references must exist either in
  persisted state or in earlier ops within the same batch.
\item Idempotency check: if the \texttt{opBatchId} has been seen before,
  broadcast \texttt{OpsAccepted} with the known revision range.
\item Lock conflict check: for operations requiring notation locks
  (\texttt{UpdateViewObjectOpaque}, \texttt{UpdateConnectionOpaque}),
  verify no other session holds a conflicting lock.
\item Revision range assignment: allocate \texttt{from..to} consuming one
  revision per op.
\item Causal metadata normalization: stamp each op with
  \texttt{clientId}, \texttt{lamport}, \texttt{opId}.
\item Idempotency recording.
\item Persistence: append to op-log, apply to materialized state, update
  head revision.
\item Optional post-commit consistency checks.
\item Kafka publish (with fallback direct broadcast on failure).
\item Broadcast \texttt{OpsAccepted} to all sessions.
\end{enumerate}

If any step fails before persistence, an \texttt{Error} message is broadcast
with a specific code (\texttt{REVISION\_AHEAD}, \texttt{PRECONDITION\_FAILED},
\texttt{LOCK\_CONFLICT}).

\subsubsection{Outbound messages (server to client)}

\begin{longtable}{p{36mm}p{52mm}p{66mm}}
\toprule
Type & Purpose & Key payload fields \\
\midrule
\texttt{CheckoutSnapshot} & Full model state &
  \texttt{modelId}, \texttt{headRevision},
  \texttt{elements}, \texttt{relationships},
  \texttt{views}, \texttt{folders}, etc. \\
\texttt{CheckoutDelta} & Incremental changes &
  \texttt{modelId},
  \texttt{fromRevision}, \texttt{toRevision},
  \texttt{ops} (array) \\
\texttt{OpsAccepted} & Batch commit confirmation &
  \texttt{opBatchId}, \texttt{baseRevision},
  \texttt{assignedRevisionRange} \\
\texttt{OpsBroadcast} & Remote batch broadcast &
  \texttt{opBatch} (full committed batch) \\
\texttt{LockEvent} & Lock state change &
  \texttt{action} (acquired/released/expired),
  \texttt{actor}, \texttt{targets} \\
\texttt{PresenceBroadcast} & Remote cursor/selection &
  \texttt{actor}, \texttt{viewId},
  \texttt{selection}, \texttt{cursor} \\
\texttt{Error} & Rejection or failure &
  \texttt{code}, \texttt{message} \\
\bottomrule
\end{longtable}

\subsection{REST API}

The server exposes REST endpoints under the base path \texttt{/}.

\subsubsection{Model state endpoints}

\begin{longtable}{p{54mm}p{44mm}p{56mm}}
\toprule
Endpoint & Method & Purpose \\
\midrule
\texttt{/models/\{modelId\}/snapshot} & GET & Get the current materialized
  snapshot of a model. Optional query parameter \texttt{?ref=<tagName>}. \\
\texttt{/models/\{modelId\}/rebuild} & POST & Rebuild materialized state
  from the op-log. Clears materialized nodes and replays all commits. \\
\bottomrule
\end{longtable}

Example:

\begin{lstlisting}[language=bash,caption={Get snapshot}]
curl -s http://localhost:8081/models/demo/snapshot | jq .
\end{lstlisting}

\subsubsection{Admin endpoints}

\begin{longtable}{p{54mm}p{44mm}p{56mm}}
\toprule
Endpoint & Method & Purpose \\
\midrule
\texttt{/admin/models} & GET & List registered models (catalog). \\
\texttt{/admin/models/\{modelId\}} & POST & Register a new model. \\
& DELETE & Delete a model. Query param
  \texttt{?force=true} to override active-session guard. \\
\texttt{/admin/models/\{modelId\}/status} & GET & Aggregated status:
  snapshot counts, revision, consistency. \\
\texttt{/admin/models/\{modelId\}/rebuild-and-status} & POST & Rebuild and
  return status in one call. \\
\texttt{/admin/models/\{modelId\}/window?limit=25} & GET & Model window:
  status + recent activity + recent op batches. \\
\texttt{/admin/overview?limit=25} & GET & Overview for all known
  active/recent models. \\
\texttt{/admin/models/\{modelId\}/integrity} & GET & Integrity report:
  missing references, orphan nodes. \\
\texttt{/admin/models/\{modelId\}/acl} & GET, PUT & Read or update
  per-model ACL (admin, writer, reader user lists). \\
\texttt{/admin/models/\{modelId\}/tags} & POST & Create an immutable tag
  at the current HEAD. \\
\texttt{/admin/models/\{modelId\}/export} & GET & Export model package
  (snapshot + op-log + tags + ACL). \\
\texttt{/admin/models/import?overwrite=true|false} & POST & Import a model
  package. \\
\texttt{/admin/models/\{modelId\}/compact} & POST & Run compaction on
  the model (prune old op-log entries and merge metadata). \\
\texttt{/admin/auth/diagnostics} & GET & Authorization diagnostics:
  resolved identity, subject, normalized roles. \\
\texttt{/admin/sessions} & GET & Active WebSocket sessions across all
  models. \\
\texttt{/admin/audit/config} & GET & Current audit configuration. \\
\bottomrule
\end{longtable}

Examples:

\begin{lstlisting}[language=bash,caption={Admin API examples}]
# List all models
curl -s http://localhost:8081/admin/models | jq .

# Create a model
curl -s -X POST http://localhost:8081/admin/models/demo | jq .

# Get model status
curl -s http://localhost:8081/admin/models/demo/status | jq .

# Get model window
curl -s http://localhost:8081/admin/models/demo/window?limit=25 | jq .

# Get overview
curl -s http://localhost:8081/admin/overview?limit=25 | jq .

# Delete model
curl -s -X DELETE http://localhost:8081/admin/models/demo | jq .

# Force delete (override active-session guard)
curl -s -X DELETE "http://localhost:8081/admin/models/demo?force=true" | jq .
\end{lstlisting}

\subsubsection{Admin UI endpoint}

\begin{description}
\item[Path] \texttt{/admin-ui/}
\item[Serves] Static Svelte admin application (built from
  \filepath{admin-web/}).
\end{description}

When running in dev mode, the admin UI is served from
\filepath{server/src/main/resources/META-INF/resources/admin-ui/} (populated
by the build script). In production, the Quarkus application serves the
built static assets.
